\section*{Chapter \ref{ch:b2:solid-liquid-diagrams} --- Binary Solid--Liquid Diagrams}

\begin{solution}{exo:b2:solid-liquid-diagrams:1}
At \qty{1300}{\degreeCelsius} the \omterm{def:b2:solid-liquid-diagrams:liquidus}{liquidus} is at 0.541 and the \omterm{def:b2:solid-liquid-diagrams:liquidus}{solidus} at 0.609.
0.40: liquid only. 0.58: liquid (0.541) and \omterm{def:b2:solid-liquid-diagrams:solid-solution}{solid solution} (0.609). 0.70:
\omterm{def:b2:solid-liquid-diagrams:solid-solution}{solid solution} only.
\end{solution}

\begin{solution}{exo:b2:solid-liquid-diagrams:2}
Solid fraction $(0.58 - 0.541)/(0.609 - 0.541) = 0.57$: \qty{1.15}{mol} of solid,
\qty{0.85}{mol} of liquid. Nickel: $1.15 \times 0.609 = \qty{0.70}{mol}$ in the
solid, $0.85 \times 0.541 = \qty{0.46}{mol}$ in the liquid; total \qty{1.16}{mol}
$= 2.00 \times 0.58$.
\end{solution}

\begin{solution}{exo:b2:solid-liquid-diagrams:3}
At fixed pressure $v = n + 1 - \varphi$ (no reaction). (a) $1 + 1 - 2 = 0$.
(b) $2 + 1 - 2 = 1$. (c) Liquid and two \omterm{def:b2:solid-liquid-diagrams:solid-solution}{solid solutions}: $2 + 1 - 3 = 0$.
\end{solution}

\begin{solution}{exo:b2:solid-liquid-diagrams:4}
A smooth fall to \qty{232}{\degreeCelsius}, a horizontal plateau while the tin
freezes, then a smooth fall again, slower near room temperature. On the
plateau liquid and solid coexist at fixed pressure, \omterm{def:b2:equilibrium-shifts:variance}{variance} 0: the heat
removed is supplied by the crystallisation, and the temperature cannot change
until the last liquid has frozen.
\end{solution}

\begin{solution}{exo:b2:solid-liquid-diagrams:5}
Benzene: $\ln x_B = -(9870/8.314)(1/269.6 - 1/278.64) = -0.1429$, $x_B =
0.867$. Naphthalene: $\ln x_N = -(19\,100/8.314)(1/269.6 - 1/353.2) = -2.017$,
$x_N = 0.133$. The sum is 1.000: \qty{269.6}{K} ($\qty{-3.5}{\degreeCelsius}$) is
the eutectic temperature, and the eutectic liquid has 0.133 naphthalene.
\end{solution}

\begin{solution}{exo:b2:solid-liquid-diagrams:6}
The liquid cools to the lead-side \omterm{def:b2:solid-liquid-diagrams:liquidus}{liquidus}, where crystals of (Pb) appear; the
liquid grows richer in tin, down to the eutectic at \qty{182.2}{\degreeCelsius}
and 62.1~\%. There the remaining liquid freezes into the \omterm{def:b2:solid-liquid-diagrams:eutectic}{eutectic mixture}.
Mass fraction of eutectic: $(30 - 18.91)/(62.13 - 18.91) = 0.26$. On further
cooling the solubility of tin in (Pb) decreases (dashed line): a little (Sn)
precipitates inside the (Pb) crystals.
\end{solution}

\begin{solution}{exo:b2:solid-liquid-diagrams:7}
23.3~\% NaCl by mass: $23.3/76.7 = \qty{0.304}{kg}$ of salt per kilogram of
water. At \qty{-25}{\degreeCelsius}, below the eutectic, no liquid can exist
whatever the proportions: salt and ice stay side by side as solids.
\end{solution}

\begin{solution}{exo:b2:solid-liquid-diagrams:8}
$x_{\mathrm B} = 0.5$ lies between E$_1$ and the compound: $\mathrm{AB_2}$
crystallises first, the liquid moves towards E$_1$, freezes there into the
eutectic of A and $\mathrm{AB_2}$; at the end, A + $\mathrm{AB_2}$.
$x_{\mathrm B} = 0.9$ lies between E$_2$ and B: B crystallises first, the liquid
reaches E$_2$; at the end, $\mathrm{AB_2}$ + B.
\end{solution}

\begin{solution}{exo:b2:solid-liquid-diagrams:9}
At the rear edge of the zone the liquid freezes into a solid holding only 0.78
times its copper mole fraction: the copper rejected stays in the liquid, which
moves on and grows richer, so the copper is swept towards the end of the bar.
With a ratio close to 1 each pass removes only a fraction of the impurity:
many passes are needed, each starting from the bar left by the previous one.
\end{solution}

\begin{solution}{exo:b2:solid-liquid-diagrams:10}
Derivation as in the chapter. For a dilute liquid, $\ln x_A \approx -x_B$, $1/T -
1/T^* \approx -\Delta T/T^{*2}$, so $x_B = \Delta_{\text{fus}}H_A\Delta T/(RT^{*2})$;
with $x_B \approx n_B/n_A = bM_A$, $\Delta T = (RT^{*2}M_A/\Delta_{\text{fus}}H_A)\,b$.
Benzene: $K_f = 8.314 \times 278.64^2 \times 0.07811/9870 = \qty{5.11}{K.kg/mol}$.
\end{solution}

\begin{solution}{exo:b2:solid-liquid-diagrams:11}
Per \qty{100}{g}: $19.0/24.31 = \qty{0.782}{mol}$ of Mg, $81.0/207.2 =
\qty{0.391}{mol}$ of Pb; ratio $2.00$: $\ce{Mg2Pb}$.
\end{solution}

\begin{solution}{exo:b2:solid-liquid-diagrams:12}
The plateau is proportional to the mass of eutectic liquid: the unknown holds
$120/300 = 0.40$ of it. On the lead side, $w = 18.91 + 0.40 \times (62.13 -
18.91) = 36.2~\%$ tin; on the tin side, $w = 96.59 - 0.40 \times (96.59 - 62.13) =
82.8~\%$. The first break of the \omterm{def:b2:solid-liquid-diagrams:cooling-curve}{cooling curve} (higher for the 36~\% alloy,
whose \omterm{def:b2:solid-liquid-diagrams:liquidus}{liquidus} is on the steep lead side) or a micrograph (primary (Pb) or
primary (Sn) crystals) tells them apart.
\end{solution}

\begin{solution}{pb:b2:solid-liquid-diagrams:1}
\textbf{1.} Lead \qty{327.5}{\degreeCelsius}, tin \qty{231.9}{\degreeCelsius}.
\textbf{2.} Per \qty{100}{g}: $62.13/118.71 = \qty{0.5234}{mol}$ of tin and
$37.87/207.2 = \qty{0.1828}{mol}$ of lead; $x_{\mathrm{Sn}} = 0.5234/0.7062 =
0.741$.
\textbf{3.} Points (0, 327.5), (100, 231.9), (62.13, 182.2), (18.91, 182.2),
(96.59, 182.2).
\textbf{4.} As in the figure: liquid; L + (Pb); L + (Sn); (Pb); (Sn); (Pb) + (Sn).
\textbf{5.} L (62.1~\% Sn) $\longrightarrow$ (Pb) (18.9~\%) + (Sn) (96.6~\%).
\textbf{6.} $v = 2 + 1 - 3 = 0$: the temperature and the three compositions are
fixed.
\textbf{7.} The lead-rich \omterm{def:b2:solid-liquid-diagrams:solid-solution}{solid solution} (Pb): lead dissolves tin, so the solid in
equilibrium with the liquid already holds tin, in the proportion read on the
\omterm{def:b2:solid-liquid-diagrams:liquidus}{solidus}.
\textbf{8.} It follows the \omterm{def:b2:solid-liquid-diagrams:liquidus}{liquidus} towards the eutectic, growing richer in tin,
up to 62.13~\%.
\textbf{9.} Liquid of 62.13~\% and (Pb) of 18.91~\% tin.
\textbf{10.} (Pb): $(62.13 - 40)/(62.13 - 18.91) = 0.512$; liquid 0.488.
\textbf{11.} (Pb) of 18.91~\% and (Sn) of 96.59~\%; (Sn): $(40 - 18.91)/(96.59 -
18.91) = 0.271$; (Pb): 0.729.
\textbf{12.} Eutectic constituent 48.8~\% (the liquid present just above);
primary (Pb) 51.2~\%. The (Pb) of question 11 is the sum of the primary
crystals and of the (Pb) lamellae of the eutectic.
\textbf{13.} A break at the \omterm{def:b2:solid-liquid-diagrams:liquidus}{liquidus}, a slower fall, then a plateau at
\qty{182.2}{\degreeCelsius} lasting 0.488 times that of the eutectic alloy.
\textbf{14.} $\ln x_{\mathrm{Sn}} = -(7030/8.314)(1/T - 1/505.1)$.
\textbf{15.} $\ln x_{\mathrm{Bi}} = -(11\,300/8.314)(1/T - 1/544.6)$.
\textbf{16.} $x_{\mathrm{Sn}}(T) + x_{\mathrm{Bi}}(T) = 1$.
\textbf{17.} At \qty{110}{\degreeCelsius}: $0.587 + 0.349 = 0.936$; at
\qty{120}{\degreeCelsius}: $0.621 + 0.382 = 1.003$; at
\qty{130}{\degreeCelsius}: $0.655 + 0.417 = 1.072$. The sum reaches 1 near
\qty{119.5}{\degreeCelsius}: \qty{120}{\degreeCelsius} to the nearest degree.
\textbf{18.} $x_{\mathrm{Bi}} = 0.38$; in mass $0.38 \times 208.98/(0.38 \times
208.98 + 0.62 \times 118.71) = 0.52$.
\textbf{19.} The model gives \qty{120}{\degreeCelsius} and 52~\% bismuth, the
assessed diagram \qty{138.8}{\degreeCelsius} and 57~\%: about
\qty{19}{\degreeCelsius} too low.
\textbf{20.} The tin that crystallises is not pure but a solution holding
bismuth: its \omterm{def:b2:chemical-potential:chemical-potential}{chemical potential} is lower than that of pure tin, the solid is
more stable, and it crystallises from the liquid at a higher temperature for a
given liquid composition. The tin branch of the \omterm{def:b2:solid-liquid-diagrams:liquidus}{liquidus} is raised, and it
meets the bismuth branch higher (a non-ideal liquid adds its own correction).
\textbf{21.} It melts and freezes at one temperature, the lowest of the system:
no pasty range during which a joint moved would crack.
\textbf{22.} Lead is toxic, and electronic waste ends up in the environment.
\textbf{23.} Heat-sensitive components can be soldered at a lower temperature;
but the joint softens at a lower temperature in service, and bismuth makes it
brittle.
\textbf{24.} \qty{570}{g} of bismuth and \qty{430}{g} of tin.
\textbf{25.} The ideal-liquid model predicts a eutectic at $\boldsymbol{T_E
\approx \qty{120}{\degreeCelsius}}$, against \qty{138.8}{\degreeCelsius}
assessed.
\end{solution}
